% Compile beside math-review.sty; replace all visible insertion text for publication.
% Rename and regroup functional headings around the actual mathematical thread.
% Shared files do not determine section hierarchy; no environment/count quota applies.
% Edition check: identify what each substantial addition helps the reader understand.
% Copyedit: state a full result once locally; check the printed meaning of notation.
\documentclass[12pt,reqno]{amsart}
\usepackage{math-review}
\title[Part I: Foundations (standard)]{[Review topic]\\
Part I: Foundations and Models (standard)}
\author{[Author name]}
\date{[Prepared date; literature cutoff date]}
\begin{document}
\begin{abstract}
\placeholder{Identify the central problem, model structures and conceptual relationships developed here, including the deeper foundational explanation that supports the general theory.}
\end{abstract}
\maketitle
\tableofcontents

\section{Introduction}\label{sec:introduction}
\placeholder{Explain the motivating problem and the larger structural questions that organize these foundations. Give relevant historical context only when it clarifies those relationships.}
\placeholder{Explain how Section~\ref{sec:language} prepares the model reasoning in Section~\ref{sec:models}, and how Section~\ref{sec:general} extracts its general significance. State which conceptual relationships receive deeper treatment than a concise account.}

\section{Objects and formulations of the question}\label{sec:language}
\placeholder{Develop the objects, defining properties and conventions actually needed. Explain what each definition enables in the models.}

\subsection{Formulations that change the viewpoint}\label{subsec:formulations}
\placeholder{Compare useful equivalent or neighboring formulations with the necessary justification and qualifications. Develop a preliminary argument only when it supports a later step; a separate foundations course is not required.}

\section{Core models and their structural explanation}\label{sec:models}
\placeholder{Develop the core model argument completely at the declared level. Explain why its key construction or choice works, identify dependencies and check the closing inference.}

\subsection{A revealing variation}\label{subsec:variation}
\placeholder{Where useful, change one feature to show what is essential or reusable. Deepen the same conceptual model before adding breadth; remove this subsection if no variation helps.}

\section{The general question and the role of the models}\label{sec:general}
\placeholder{Explain the exact bridge from the models to the central problem. Compare what persists and what needs a different input, without treating failure of one method as impossibility. State the general target and clarify why the developed language makes the main theory accessible.}

\templatereferences
\end{document}
